\begin{table}[tbp]
\centering
\small
\caption{Illustrative scenarios for the absolute decouplers above the boundary.}
\label{tab:scenarios}
\begin{threeparttable}
\begin{tabular}{lrrr}
\toprule
Scenario & Median cut (\%/yr) & Reach 1.61 t by 2050 & Median year reached \\
\midrule
Achieved 2011--2022 rates & 1.6 & 2 of 41 & 2115 \\
Each country doubles its rate & 3.2 & 6 of 41 & 2068 \\
All at the 90th-percentile rate & 3.2 & 7 of 41 & 2071 \\
All at the fastest achieved rate & 3.9 & 10 of 41 & 2061 \\
Each at the rate required by 2050 & 5.5 & 41 of 41 & 2050 \\
\bottomrule
\end{tabular}
\begin{tablenotes}[flushleft]
\footnotesize
\item \textit{Notes:} Each scenario applies a constant annual rate of change in consumption-based CO$_2$ per person from each country's 2022 level. The 90th-percentile and fastest rates are those achieved by this group in 2011--2022. These are arithmetic illustrations, not forecasts.
\end{tablenotes}
\end{threeparttable}
\end{table}
